\documentclass[a4paper,11pt]{article}

\usepackage[utf8]{inputenc}
\usepackage[italian]{babel}
\usepackage{geometry}
\usepackage{enumitem}
\usepackage{booktabs}
\usepackage{amsmath}
\usepackage{amssymb}
\usepackage{xcolor}
\usepackage{graphicx}
\usepackage{tikz}
\usepackage{float}
\usepackage{multirow}
\usepackage{fancyhdr}
\usepackage{hyperref}

\usetikzlibrary{shapes.geometric, arrows.meta, positioning, fit, backgrounds}

\geometry{a4paper, margin=2.5cm}

\hypersetup{
    colorlinks=true,
    linkcolor=blue!60!black,
    urlcolor=blue!60!black
}

\pagestyle{fancy}
\fancyhf{}
\rhead{SmartLight Pisa - Shadow Mode}
\lhead{Documento Tecnico}
\rfoot{Pagina \thepage}

\definecolor{completed}{RGB}{34, 139, 34}
\definecolor{pending}{RGB}{200, 80, 80}
\definecolor{partial}{RGB}{230, 150, 50}

\title{
    \vspace{-1cm}
    \textbf{SmartLight Pisa}\\
    \Large Piano di Testing in Shadow Mode\\
    \large Documento Tecnico e Analisi dei Costi
}
\author{SmartLight Research Team}
\date{29 Gennaio 2026}

\begin{document}

\maketitle

\begin{abstract}
Questo documento presenta il piano dettagliato per l'implementazione del primo test in \textbf{Shadow Mode} del sistema SmartLight Pisa. Il Shadow Mode consente di validare le decisioni dell'agente AI in ambiente reale senza controllo fisico dei semafori, minimizzando i rischi e massimizzando la raccolta di dati. Il documento include: stato attuale del progetto, architettura del sistema di test con edge computing remoto, analisi dei costi con intervalli di confidenza ($\mu \pm \sigma$) e timeline di implementazione.
\end{abstract}

\tableofcontents
\newpage

%==============================================================================
\section{Executive Summary: Stato del Progetto}
%==============================================================================

\subsection{Progresso Complessivo}

Il progetto SmartLight Pisa ha raggiunto la versione \textbf{v3.4}, completando con successo le fasi fondamentali di ottimizzazione e accelerazione hardware. La tabella seguente riassume lo stato attuale:

\begin{table}[H]
    \centering
    \begin{tabular}{clcc}
        \toprule
        \textbf{Fase} & \textbf{Descrizione} & \textbf{Stato} & \textbf{Completamento} \\
        \midrule
        1 & Training Optimization & \textcolor{completed}{\textbf{Completata}} & 100\% \\
        2 & Advanced Algorithms & \textcolor{partial}{\textbf{Parziale}} & 66\% \\
        3 & Hardware Acceleration & \textcolor{completed}{\textbf{Completata}} & 100\% \\
        4 & Multi-Agent Coordination & \textcolor{pending}{\textbf{Pending}} & 0\% \\
        5 & Real-World Perception & \textcolor{pending}{\textbf{Pending}} & 0\% \\
        6 & Robustness \& Deployment & \textcolor{pending}{\textbf{Pending}} & 0\% \\
        \bottomrule
    \end{tabular}
    \caption{Stato delle fasi di sviluppo al 29 Gennaio 2026}
\end{table}

\subsection{Feature Implementate (v3.4)}

L'attuale versione del sistema include le seguenti ottimizzazioni:

\begin{enumerate}[label=\arabic*.]
    \item \textbf{Prioritized Experience Replay} - Campionamento pesato basato su TD-error:
    \begin{equation}
        P(i) = \frac{p_i^\alpha}{\sum_k p_k^\alpha}, \quad p_i = |\delta_i| + \epsilon
    \end{equation}
    
    \item \textbf{Learning Rate Scheduler} - Decay adattivo con StepLR ($\gamma = 0.7$, step $= 50$)
    
    \item \textbf{State Compression} - Riduzione dello spazio degli stati da 17 a 11 dimensioni
    
    \item \textbf{Batch Normalization} - Stabilizzazione del training con normalizzazione per layer
    
    \item \textbf{Dueling DQN Architecture} - Separazione dei flussi Value e Advantage:
    \begin{equation}
        Q(s, a) = V(s) + \left( A(s, a) - \frac{1}{|\mathcal{A}|} \sum_{a'} A(s, a') \right)
    \end{equation}
    
    \item \textbf{Double DQN} - Riduzione dell'overestimation dei Q-values:
    \begin{equation}
        Y_t^{\text{DDQN}} = R_{t+1} + \gamma Q\left(S_{t+1}, \arg\max_a Q(S_{t+1}, a; \theta_t); \theta_t^-\right)
    \end{equation}
    
    \item \textbf{GPU Acceleration} - Supporto nativo CUDA (testato su NVIDIA GTX 1650)
    
    \item \textbf{Parallel Environment Simulation} - 4 ambienti simultanei ($\sim$200 transizioni/secondo)
\end{enumerate}

\subsection{Prossimi Step Prioritari}

\begin{table}[H]
    \centering
    \begin{tabular}{clcl}
        \toprule
        \textbf{Priorità} & \textbf{Feature} & \textbf{Fase} & \textbf{Impatto} \\
        \midrule
        \textbf{CRITICO} & YOLOv8 Integration & 5 & Bridge simulazione $\rightarrow$ realtà \\
        ALTO & Shadow Mode Testing & 5-6 & Validazione decisioni AI \\
        ALTO & Multi-Agent & 4 & Scalabilità multi-incrocio \\
        MEDIO & PPO Transition & 2 & Stabilità policy \\
        \bottomrule
    \end{tabular}
    \caption{Priorità di sviluppo}
\end{table}

%==============================================================================
\section{Shadow Mode: Definizione e Obiettivi}
%==============================================================================

\subsection{Cos'è il Shadow Mode}

Il \textbf{Shadow Mode} è un paradigma di testing in cui l'agente AI:

\begin{itemize}
    \item \textbf{Osserva} il traffico reale attraverso videocamere
    \item \textbf{Calcola} la decisione ottimale secondo il modello addestrato
    \item \textbf{NON controlla} fisicamente i semafori
    \item \textbf{Compara} la propria decisione con quella del controllore esistente
    \item \textbf{Registra} tutte le divergenze per analisi successiva
\end{itemize}

\subsection{Vantaggi del Shadow Mode}

\begin{enumerate}
    \item \textbf{Zero rischi} - Nessun impatto sul traffico reale
    \item \textbf{Validazione empirica} - Test con dati reali, non simulati
    \item \textbf{Identificazione edge cases} - Scenari non coperti dal training
    \item \textbf{Baseline comparison} - Confronto diretto AI vs. sistema attuale
    \item \textbf{Raccolta dati} - Dataset reale per fine-tuning
\end{enumerate}

\subsection{Obiettivi Misurabili}

\begin{table}[H]
    \centering
    \begin{tabular}{lcc}
        \toprule
        \textbf{Metrica} & \textbf{Target Minimo} & \textbf{Target Ottimo} \\
        \midrule
        Agreement Rate ($\pm 5s$) & $\geq 70\%$ & $\geq 85\%$ \\
        MAE (Mean Absolute Error) & $\leq 8s$ & $\leq 4s$ \\
        Emergency Response Accuracy & $\geq 95\%$ & $100\%$ \\
        Uptime del sistema & $\geq 95\%$ & $\geq 99\%$ \\
        Decisioni raccolte/giorno & $\geq 1000$ & $\geq 5000$ \\
        \bottomrule
    \end{tabular}
    \caption{KPI per il Shadow Mode Testing}
\end{table}

%==============================================================================
\section{Architettura del Sistema}
%==============================================================================

\subsection{Panoramica dell'Architettura}

L'architettura proposta prevede una separazione fisica tra i componenti di acquisizione (on-site) e quelli di elaborazione (remoto), per motivi di sicurezza, manutenibilità e costi.

\begin{figure}[H]
    \centering
    \begin{tikzpicture}[
        node distance=1cm,
        box/.style={rectangle, draw, rounded corners, minimum height=1cm, minimum width=3cm, align=center, fill=blue!10},
        smallbox/.style={rectangle, draw, rounded corners, minimum height=0.8cm, minimum width=2.5cm, align=center, fill=green!10},
        cloudbox/.style={rectangle, draw, rounded corners, minimum height=1cm, minimum width=3.5cm, align=center, fill=orange!15, dashed},
        arrow/.style={-{Stealth[scale=1.2]}, thick}
    ]
    
    % ON-SITE COMPONENTS
    \node[box] (camera) {Telecamera IP\\1080p/30fps};
    \node[box, right=2cm of camera] (encoder) {Video Encoder\\H.264/H.265};
    \node[smallbox, below=0.8cm of camera] (semaphore) {Stato Semaforo\\(API/PLC)};
    
    % NETWORK
    \node[cloudbox, right=2cm of encoder] (network) {Rete 4G/5G\\o Fibra};
    
    % REMOTE COMPONENTS
    \node[box, right=2cm of network] (xavier) {NVIDIA Jetson\\Xavier NX};
    \node[smallbox, above=0.8cm of xavier] (yolo) {YOLOv8\\Detection};
    \node[smallbox, below=0.8cm of xavier] (dqn) {Dueling DQN\\v3.4};
    
    % OUTPUT
    \node[box, right=2cm of xavier, fill=yellow!20] (comparator) {Shadow\\Comparator};
    \node[smallbox, right=1.5cm of comparator] (dashboard) {Dashboard\\\& Logs};
    
    % ARROWS
    \draw[arrow] (camera) -- node[above, font=\small] {stream} (encoder);
    \draw[arrow] (encoder) -- (network);
    \draw[arrow] (network) -- (xavier);
    \draw[arrow] (semaphore) -| node[below, font=\small, pos=0.3] {stato} (network);
    \draw[arrow] (yolo) -- (xavier);
    \draw[arrow] (xavier) -- (dqn);
    \draw[arrow] (xavier) -- (comparator);
    \draw[arrow] (comparator) -- (dashboard);
    
    % SITE BOUNDARIES
    \begin{scope}[on background layer]
        \node[draw, dashed, thick, rounded corners, fill=blue!5, 
              fit=(camera)(encoder)(semaphore), 
              label=above:\textbf{ON-SITE (Incrocio)}] {};
        \node[draw, dashed, thick, rounded corners, fill=orange!5,
              fit=(xavier)(yolo)(dqn)(comparator)(dashboard),
              label=above:\textbf{REMOTO (Data Center / Ufficio)}] {};
    \end{scope}
    
    \end{tikzpicture}
    \caption{Architettura Shadow Mode con elaborazione remota}
\end{figure}

\subsection{Componenti On-Site (Incrocio)}

I componenti installati presso l'incrocio sono minimali per ridurre costi di manutenzione e rischi di guasto:

\begin{table}[H]
    \centering
    \begin{tabular}{lll}
        \toprule
        \textbf{Componente} & \textbf{Funzione} & \textbf{Specifiche} \\
        \midrule
        Telecamera IP & Acquisizione video & 1080p, 30fps, PoE, IP66 \\
        Box Encoder & Compressione stream & H.265, basso bitrate \\
        Gateway 4G/5G & Trasmissione dati & LTE Cat.6+, IP statico \\
        Sensore Semaforo & Lettura stato fasi & Interfaccia PLC/API \\
        \bottomrule
    \end{tabular}
    \caption{Hardware on-site}
\end{table}

\subsection{Componenti Remoti (Edge Computing)}

L'elaborazione avviene su un dispositivo \textbf{NVIDIA Jetson Xavier NX} situato in posizione remota (ufficio, data center locale o cloud edge):

\begin{table}[H]
    \centering
    \begin{tabular}{ll}
        \toprule
        \textbf{Specifica} & \textbf{Valore} \\
        \midrule
        GPU & 384 CUDA cores + 48 Tensor cores \\
        CPU & 6-core NVIDIA Carmel ARM v8.2 \\
        Memoria & 8 GB 128-bit LPDDR4x \\
        Storage & 256 GB NVMe SSD \\
        AI Performance & 21 TOPS \\
        Potenza & $10-20W$ (modalità configurabili) \\
        \bottomrule
    \end{tabular}
    \caption{Specifiche NVIDIA Jetson Xavier NX}
\end{table}

\subsubsection{Perché Xavier Remoto?}

\begin{enumerate}
    \item \textbf{Sicurezza} - Nessun hardware critico esposto a intemperie/vandalismi
    \item \textbf{Manutenibilità} - Accesso fisico facile per aggiornamenti/debug
    \item \textbf{Scalabilità} - Un Xavier può gestire più incroci simultaneamente
    \item \textbf{Costi} - Connettività 4G/5G più economica di installazione IP industriale
    \item \textbf{Latenza accettabile} - Per shadow mode non serve real-time stretto ($<100ms$ sufficiente)
\end{enumerate}

\subsection{Flusso Dati}

\begin{enumerate}
    \item \textbf{Acquisizione} ($t_0$): Telecamera cattura frame 1080p@30fps
    \item \textbf{Compressione} ($t_0 + 10ms$): Encoder comprime in H.265 (2-4 Mbps)
    \item \textbf{Trasmissione} ($t_0 + 50ms$): Stream inviato via 4G/5G al server remoto
    \item \textbf{Decodifica} ($t_0 + 70ms$): Xavier decodifica frame
    \item \textbf{Detection} ($t_0 + 100ms$): YOLOv8 rileva veicoli e pedoni
    \item \textbf{Inference} ($t_0 + 110ms$): DQN calcola azione ottimale
    \item \textbf{Comparazione} ($t_0 + 115ms$): Confronto con decisione reale
    \item \textbf{Logging} ($t_0 + 120ms$): Salvataggio su database
\end{enumerate}

\textbf{Latenza end-to-end stimata}: $100-150ms$ (accettabile per shadow mode).

%==============================================================================
\section{Analisi dei Costi}
%==============================================================================

In questa sezione presentiamo un'analisi dettagliata dei costi, con intervalli di confidenza basati sulla qualità del test. I costi sono espressi come $\mu \pm \sigma$, dove:

\begin{itemize}
    \item $\mu$ = costo medio atteso
    \item $\sigma$ = deviazione standard, funzione della qualità/affidabilità
\end{itemize}

La varianza $\sigma$ è modellata come:
\begin{equation}
    \sigma = k \cdot \mu \cdot (1 - Q)
\end{equation}
dove $Q \in [0, 1]$ è l'indice di qualità del test (0 = minimo, 1 = massimo) e $k$ è un fattore di scala tipicamente $\in [0.1, 0.3]$.

\subsection{Costi Hardware}

\subsubsection{Scenario A: Test Minimo (Proof of Concept)}

Qualità test: $Q = 0.4$, $k = 0.25$

\begin{table}[H]
    \centering
    \begin{tabular}{lrrr}
        \toprule
        \textbf{Componente} & \textbf{$\mu$ (€)} & \textbf{$\sigma$ (€)} & \textbf{Range (€)} \\
        \midrule
        Webcam USB 1080p & 50 & 7.50 & $42.50 - 57.50$ \\
        Mini PC (usato) & 150 & 22.50 & $127.50 - 172.50$ \\
        Chiavetta 4G & 30 & 4.50 & $25.50 - 34.50$ \\
        Cavi e accessori & 30 & 4.50 & $25.50 - 34.50$ \\
        \midrule
        \textbf{Subtotale HW} & \textbf{260} & \textbf{39} & $\mathbf{221 - 299}$ \\
        \bottomrule
    \end{tabular}
    \caption{Costi hardware - Scenario Minimo ($Q = 0.4$)}
\end{table}

\subsubsection{Scenario B: Test Standard (Raccomandato)}

Qualità test: $Q = 0.7$, $k = 0.20$

\begin{table}[H]
    \centering
    \begin{tabular}{lrrr}
        \toprule
        \textbf{Componente} & \textbf{$\mu$ (€)} & \textbf{$\sigma$ (€)} & \textbf{Range (€)} \\
        \midrule
        Telecamera IP PoE (Hikvision) & 180 & 10.80 & $169.20 - 190.80$ \\
        NVIDIA Jetson Xavier NX & 450 & 27.00 & $423.00 - 477.00$ \\
        Router 4G industriale & 200 & 12.00 & $188.00 - 212.00$ \\
        SSD NVMe 256GB & 50 & 3.00 & $47.00 - 53.00$ \\
        Alimentatore PoE & 40 & 2.40 & $37.60 - 42.40$ \\
        Custodia IP66 & 80 & 4.80 & $75.20 - 84.80$ \\
        Cavi e accessori & 50 & 3.00 & $47.00 - 53.00$ \\
        \midrule
        \textbf{Subtotale HW} & \textbf{1050} & \textbf{63} & $\mathbf{987 - 1113}$ \\
        \bottomrule
    \end{tabular}
    \caption{Costi hardware - Scenario Standard ($Q = 0.7$)}
\end{table}

\subsubsection{Scenario C: Test Professionale}

Qualità test: $Q = 0.9$, $k = 0.15$

\begin{table}[H]
    \centering
    \begin{tabular}{lrrr}
        \toprule
        \textbf{Componente} & \textbf{$\mu$ (€)} & \textbf{$\sigma$ (€)} & \textbf{Range (€)} \\
        \midrule
        Telecamera PTZ IP (Axis) & 600 & 9.00 & $591.00 - 609.00$ \\
        NVIDIA Jetson AGX Orin & 1500 & 22.50 & $1477.50 - 1522.50$ \\
        Router 5G industriale & 400 & 6.00 & $394.00 - 406.00$ \\
        SSD NVMe 512GB & 80 & 1.20 & $78.80 - 81.20$ \\
        UPS compatto & 150 & 2.25 & $147.75 - 152.25$ \\
        Custodia IP67 ventilata & 200 & 3.00 & $197.00 - 203.00$ \\
        Sistema di mounting & 100 & 1.50 & $98.50 - 101.50$ \\
        Cavi industriali & 70 & 1.05 & $68.95 - 71.05$ \\
        \midrule
        \textbf{Subtotale HW} & \textbf{3100} & \textbf{46.50} & $\mathbf{3053.50 - 3146.50}$ \\
        \bottomrule
    \end{tabular}
    \caption{Costi hardware - Scenario Professionale ($Q = 0.9$)}
\end{table}

\subsection{Costi Operativi (Mensili)}

I costi operativi sono ricorrenti e dipendono dalla durata del test.

\begin{table}[H]
    \centering
    \begin{tabular}{lrrrr}
        \toprule
        \textbf{Voce} & \textbf{Scenario A} & \textbf{Scenario B} & \textbf{Scenario C} & \textbf{$\sigma$} \\
        \midrule
        SIM dati 4G/5G (50GB) & €15 & €25 & €50 & $\pm 5$ \\
        Elettricità (kWh) & €5 & €10 & €20 & $\pm 2$ \\
        Cloud storage (100GB) & €5 & €10 & €25 & $\pm 3$ \\
        Manutenzione remota & €0 & €50 & €150 & $\pm 20$ \\
        \midrule
        \textbf{Totale mensile} & \textbf{€25} & \textbf{€95} & \textbf{€245} & $\pm 30$ \\
        \bottomrule
    \end{tabular}
    \caption{Costi operativi mensili per scenario}
\end{table}

\subsection{Costi di Sviluppo Software}

\begin{table}[H]
    \centering
    \begin{tabular}{lrrr}
        \toprule
        \textbf{Attività} & \textbf{Ore Stimate} & \textbf{Costo/Ora (€)} & \textbf{Totale (€)} \\
        \midrule
        Integrazione YOLOv8 & $20 \pm 5$ & 50 & $1000 \pm 250$ \\
        Script Shadow Mode & $15 \pm 3$ & 50 & $750 \pm 150$ \\
        Dashboard monitoring & $25 \pm 8$ & 50 & $1250 \pm 400$ \\
        Testing e debug & $20 \pm 10$ & 50 & $1000 \pm 500$ \\
        Documentazione & $10 \pm 2$ & 50 & $500 \pm 100$ \\
        \midrule
        \textbf{Totale sviluppo} & $\mathbf{90 \pm 28}$ & -- & $\mathbf{4500 \pm 1400}$ \\
        \bottomrule
    \end{tabular}
    \caption{Costi di sviluppo software}
\end{table}

\subsection{Riepilogo Costi Totali}

Assumendo un test di \textbf{3 mesi} di durata:

\begin{table}[H]
    \centering
    \begin{tabular}{lccc}
        \toprule
        \textbf{Categoria} & \textbf{Scenario A (€)} & \textbf{Scenario B (€)} & \textbf{Scenario C (€)} \\
        \midrule
        Hardware (una tantum) & $260 \pm 39$ & $1050 \pm 63$ & $3100 \pm 47$ \\
        Operativi (3 mesi) & $75 \pm 15$ & $285 \pm 45$ & $735 \pm 60$ \\
        Sviluppo software & \multicolumn{3}{c}{$4500 \pm 1400$} \\
        \midrule
        \textbf{TOTALE} & $\mathbf{4835 \pm 1454}$ & $\mathbf{5835 \pm 1508}$ & $\mathbf{8335 \pm 1507}$ \\
        \midrule
        \textbf{Range IC 95\%} & €1927 - €7743 & €2819 - €8851 & €5321 - €11349 \\
        \bottomrule
    \end{tabular}
    \caption{Riepilogo costi totali per 3 mesi di test}
\end{table}

\textbf{Nota}: L'intervallo di confidenza al 95\% è calcolato come $\mu \pm 2\sigma$.

\subsection{Analisi Costo-Beneficio}

Definiamo il \textbf{Value Index} come:
\begin{equation}
    VI = \frac{Q \times D \times R}{C_{tot}}
\end{equation}
dove:
\begin{itemize}
    \item $Q$ = Qualità del test (0-1)
    \item $D$ = Numero di decisioni raccolte (migliaia)
    \item $R$ = Affidabilità del sistema (uptime \%)
    \item $C_{tot}$ = Costo totale in migliaia di euro
\end{itemize}

\begin{table}[H]
    \centering
    \begin{tabular}{lccccc}
        \toprule
        \textbf{Scenario} & $Q$ & $D$ (k) & $R$ (\%) & $C_{tot}$ (k€) & \textbf{VI} \\
        \midrule
        A (Minimo) & 0.4 & 90 & 85 & 4.84 & 6.3 \\
        B (Standard) & 0.7 & 270 & 95 & 5.84 & 30.8 \\
        C (Professionale) & 0.9 & 450 & 99 & 8.34 & 48.1 \\
        \bottomrule
    \end{tabular}
    \caption{Analisi Value Index per scenario}
\end{table}

\textbf{Raccomandazione}: Lo \textbf{Scenario B} offre il miglior compromesso costo/qualità per un primo test pilota.

%==============================================================================
\section{Timeline di Implementazione}
%==============================================================================

\subsection{Gantt del Progetto}

\begin{table}[H]
    \centering
    \begin{tabular}{lccccc}
        \toprule
        \textbf{Attività} & \textbf{Sett 1-2} & \textbf{Sett 3-4} & \textbf{Sett 5-8} & \textbf{Sett 9-12} & \textbf{Sett 13+} \\
        \midrule
        Approvvigionamento HW & $\bullet\bullet$ & & & & \\
        Sviluppo YOLOv8 integration & $\bullet$ & $\bullet\bullet$ & & & \\
        Setup telecamera test & & $\bullet\bullet$ & & & \\
        Sviluppo Shadow Mode script & & $\bullet$ & $\bullet\bullet$ & & \\
        Testing interno & & & $\bullet\bullet$ & & \\
        Deployment on-site & & & & $\bullet$ & \\
        Raccolta dati Shadow Mode & & & & $\bullet\bullet$ & $\bullet\bullet$ \\
        Analisi e reportistica & & & & & $\bullet\bullet$ \\
        \bottomrule
    \end{tabular}
    \caption{Timeline del progetto ($\bullet$ = bassa intensità, $\bullet\bullet$ = alta intensità)}
\end{table}

\subsection{Milestone Principali}

\begin{enumerate}
    \item \textbf{M1 (Settimana 2)}: Hardware ordinato e ricevuto
    \item \textbf{M2 (Settimana 4)}: YOLOv8 integrato e testato su video registrati
    \item \textbf{M3 (Settimana 6)}: Script Shadow Mode funzionante in laboratorio
    \item \textbf{M4 (Settimana 8)}: Sistema completo testato end-to-end
    \item \textbf{M5 (Settimana 10)}: Deployment presso incrocio pilota
    \item \textbf{M6 (Settimana 14)}: Prima analisi dati (1000+ decisioni)
    \item \textbf{M7 (Settimana 20)}: Report finale con raccomandazioni
\end{enumerate}

%==============================================================================
\section{Metriche e KPI}
%==============================================================================

\subsection{Metriche Primarie}

\begin{enumerate}
    \item \textbf{Agreement Rate} (AR):
    \begin{equation}
        AR = \frac{\#\{|t_{AI} - t_{actual}| \leq \Delta\}}{N_{total}} \times 100\%
    \end{equation}
    dove $\Delta = 5s$ è la tolleranza e $N_{total}$ è il numero totale di decisioni.
    
    \item \textbf{Mean Absolute Error} (MAE):
    \begin{equation}
        MAE = \frac{1}{N} \sum_{i=1}^{N} |t_{AI}^{(i)} - t_{actual}^{(i)}|
    \end{equation}
    
    \item \textbf{Emergency Response Score} (ERS):
    \begin{equation}
        ERS = \frac{\#\{\text{correct emergency responses}\}}{\#\{\text{total emergency events}\}} \times 100\%
    \end{equation}
\end{enumerate}

\subsection{Metriche Secondarie}

\begin{itemize}
    \item \textbf{System Uptime}: percentuale di tempo operativo
    \item \textbf{Inference Latency}: tempo medio di elaborazione per frame
    \item \textbf{Detection Accuracy}: precisione YOLOv8 su veicoli reali
    \item \textbf{Edge Case Rate}: frequenza di decisioni fortemente divergenti ($|t_{AI} - t_{actual}| > 15s$)
\end{itemize}

%==============================================================================
\section{Rischi e Mitigazioni}
%==============================================================================

\begin{table}[H]
    \centering
    \begin{tabular}{p{4cm}cp{4cm}p{3cm}}
        \toprule
        \textbf{Rischio} & \textbf{Probabilità} & \textbf{Impatto} & \textbf{Mitigazione} \\
        \midrule
        Connettività 4G/5G instabile & Media & Perdita dati & Buffering locale, retry automatico \\
        Illuminazione scarsa (notte) & Alta & Detection degradata & Camera IR, training notturno \\
        Occlusioni veicoli & Media & Conteggi errati & Multi-camera, tracking \\
        Guasto hardware & Bassa & Downtime & Hardware ridondante \\
        API semaforo non disponibile & Media & Nessun confronto & Log manuale iniziale \\
        \bottomrule
    \end{tabular}
    \caption{Matrice dei rischi}
\end{table}

%==============================================================================
\section{Conclusioni}
%==============================================================================

Il test in Shadow Mode rappresenta il passaggio fondamentale dalla simulazione alla validazione in ambiente reale. L'architettura proposta con \textbf{elaborazione remota su Xavier NX} offre:

\begin{itemize}
    \item \textbf{Sicurezza}: hardware critico protetto
    \item \textbf{Flessibilità}: facile manutenzione e aggiornamento
    \item \textbf{Scalabilità}: possibilità di estendere a più incroci
    \item \textbf{Costi contenuti}: investimento iniziale €5000-6000 (Scenario B)
\end{itemize}

\subsection{Raccomandazione Finale}

Si raccomanda di procedere con lo \textbf{Scenario B (Standard)} per il primo test pilota, con la seguente roadmap:

\begin{enumerate}
    \item \textbf{Fase 0 (ora)}: Approvazione budget e procurement hardware
    \item \textbf{Fase 1 (2-4 settimane)}: Sviluppo integrazione YOLOv8 + Shadow Mode
    \item \textbf{Fase 2 (4-8 settimane)}: Testing interno e ottimizzazione
    \item \textbf{Fase 3 (8-20 settimane)}: Deployment e raccolta dati reali
    \item \textbf{Fase 4 (20+ settimane)}: Analisi, re-training e preparazione deployment attivo
\end{enumerate}

\vspace{1cm}
\begin{center}
    \textbf{Budget richiesto (Scenario B, 3 mesi)}: $\mathbf{5835 \pm 1508}$ \textbf{EUR}
\end{center}

\end{document}
